% Generated draft: requires mathematical and source review.
\documentclass[11pt]{article}
\usepackage[margin=25mm]{geometry}
\usepackage{fontspec}
\setmainfont{Latin Modern Roman}
\usepackage{amsmath,amssymb,mathtools,unicode-math}
\setmathfont{Latin Modern Math}
\usepackage{enumitem,needspace,xurl,hyperref}
\setlength{\parindent}{0pt}
\setlength{\parskip}{5pt}
\setlength{\emergencystretch}{3em}
\newcommand{\N}{\mathbb{N}}
\newcommand{\Z}{\mathbb{Z}}
\newcommand{\Q}{\mathbb{Q}}
\newcommand{\R}{\mathbb{R}}
\newcommand{\C}{\mathbb{C}}
\newcommand{\F}{\mathbb{F}}
\newcommand{\A}{\mathbb{A}}
\newcommand{\PP}{\mathbb{P}}
\newcommand{\E}{\mathbb{E}}
\newcommand{\eps}{\varepsilon}
\newcommand{\dd}{\,\mathrm{d}}
\begin{document}
\section*{UnsolvedMath ID: 6700068}
\textbf{Source problem code:} catalogue entry
\subsection{Definitions and mathematical statement}
Let \(B_{\mathbb H}^n(R)\) denote a radius-\(R\) ball in real hyperbolic \(n\)-space, with the conventional curvature-\(-1\) metric. Given a specified metric target \((S,d_S)\), the usual notation for the space of maps of Lipschitz constant at most one is
\[
 \operatorname{Lip}_1(B_{\mathbb H}^n(R),S)=
 \{f:d_S(f(x),f(y))\le d_{\mathbb H}(x,y)\text{ for all }x,y\in B_{\mathbb H}^n(R)\}.
\]
The source's part [b] asks for the asymptotic behavior, as \(R\to\infty\), of
\[
 D(R)=\operatorname{diam}_c\bigl(\operatorname{Lip}_1(B_{\mathbb H}^n(R),S)\bigr).
\]
The target \(S\) and the subscripted diameter \(\operatorname{diam}_c\) are not defined in the excerpt, so this is a retained source expression rather than a numerical asymptotic conjecture.

In particular the map-space metric, its intrinsic versus extrinsic nature, whether components or homotopies are involved, and any boundary constraints must be specified. Ordinary uniform distance, intrinsic path distance within the Lipschitz subspace, and a componentwise or chain diameter need not give the same function. No one is silently substituted for \(\operatorname{diam}_c\). The ball's open/closed convention, dimension range, and precise target geometry also need the omitted preceding context. The notation \(\operatorname{Lip}_1\) is read as ``at most one'', not exactly one.

The supplied summary calls this a specific open spaces-of-maps asymptotic question and reports no resolution located, but supplies no additional definitions or theorem references. That review report and the 2017 open-item tag do not determine an independently verified current status.
\subsection{Short English statement}
How does the intended diameter of the space of 1-Lipschitz maps from a growing hyperbolic ball to \(S\) behave as the radius tends to infinity? The target and diameter convention are missing, so the asymptotic quantity needs the original context.
\subsection{Sources}
\begin{itemize}
\item Local registry: \url{data/unsolvedmath/problems.json}, permanent UnsolvedMath ID \(6700068\), source code catalogue entry; supplied statement, background, and research summary, updated 2026-08-06.
\item Misha Gromov, \emph{101 Questions, Problems and Conjectures around Scalar Curvature} (2017), printed tag [?72], occurrence 71, part [b], PDF page 72: \url{https://www.ihes.fr/~gromov/wp-content/uploads/2018/08/101-problemsOct1-2017.pdf}. Supplied normalized PDF-text extraction accessed 2026-07-29; described as incomplete.
\item No separate definition of \(S\), \(\operatorname{diam}_c\), or asymptotic resolution reference is supplied.
\end{itemize}
\end{document}
